%% kapa: exam=1 round=23 year=2012 period=2 form=A subject=accounting subject_ko=회계학 seq=30 number=30
%% source: Q-Net 감정평가사 기출문제 (cst00309) artlSeq=5052580 file=raw/1008/5052580/2012년 제23회 감정평가사 제1차시험 기출문제(2교시 A형).pdf
%% conversion: kordoc 4.24.51 -> tools/gen1.py (auto); review=proofread-sonnet; flags=none
\begin{question}{30}
(주)서울은 (주)감평과 협의하여 사용하던 기계장치(취득원가 ￦8,000,000, 감가상각누계액 ￦3,000,000, 공정가치 ￦4,500,000)를 제공하는 조건으로 차량운반구를 취득하였다. 차량운반구의 판매가격은 ￦8,000,000이며, (주)감평은 이 기계장치의 가치를 ￦5,000,000으로 인정하기로 하였으며, (주)서울은 차액 ￦3,000,000을 현금으로 지급하였다. 이 교환거래는 상업적 실질이 있고, 공정가치는 합리적으로 측정한 금액이다. 차량운반구의 취득원가는?
\begin{choices}
  \item ￦7,000,000
  \item ￦7,500,000
  \item ￦8,000,000
  \item ￦8,500,000
  \item ￦9,000,000
\end{choices}
\end{question}
\answer{2}
